\chapter{Lesson 3: Katakana}

\section{What katakana is for}

Katakana (\ja{カタカナ}) has the same sound values as hiragana but angular shapes. It
developed from \emph{pieces} of kanji, used by Buddhist monks as reading notes. Today it
is used for:
\begin{itemize}
  \item loanwords (\ja{がいらいご}, \emph{gairaigo}): \ja{コーヒー} \emph{kōhī}
        ``coffee'', \ja{テレビ} \emph{terebi} ``TV'';
  \item foreign names and places: \ja{カナダ} \emph{Kanada}, \ja{トロント} \emph{Toronto},
        \ja{ウォータールー} \emph{Wōtārū} ``Waterloo'';
  \item onomatopoeia and sound effects (especially in manga): \ja{ドキドキ}
        ``heart pounding'';
  \item emphasis, like italics in English; scientific names of plants and animals.
\end{itemize}

\section{The katakana chart}

\begin{center}
\begin{tikzpicture}[x=1.25cm,y=-1.05cm, cell/.style={draw, minimum width=1.25cm, minimum height=1.05cm, inner sep=0pt}]
  \node[font=\small\bfseries] at (1,-0.15) {a};
  \node[font=\small\bfseries] at (2,-0.15) {i};
  \node[font=\small\bfseries] at (3,-0.15) {u};
  \node[font=\small\bfseries] at (4,-0.15) {e};
  \node[font=\small\bfseries] at (5,-0.15) {o};
  \node[font=\small\bfseries, anchor=east] at (0.35,0.6) {};
  \node[cell] at (1,0.6) {};
  \node[font=\Large] at (1,0.5) {\ja{ア}};
  \node[font=\scriptsize, text=gray!50!black] at (1,0.92) {a};
  \node[cell] at (2,0.6) {};
  \node[font=\Large] at (2,0.5) {\ja{イ}};
  \node[font=\scriptsize, text=gray!50!black] at (2,0.92) {i};
  \node[cell] at (3,0.6) {};
  \node[font=\Large] at (3,0.5) {\ja{ウ}};
  \node[font=\scriptsize, text=gray!50!black] at (3,0.92) {u};
  \node[cell] at (4,0.6) {};
  \node[font=\Large] at (4,0.5) {\ja{エ}};
  \node[font=\scriptsize, text=gray!50!black] at (4,0.92) {e};
  \node[cell] at (5,0.6) {};
  \node[font=\Large] at (5,0.5) {\ja{オ}};
  \node[font=\scriptsize, text=gray!50!black] at (5,0.92) {o};
  \node[font=\small\bfseries, anchor=east] at (0.35,1.6) {k};
  \node[cell] at (1,1.6) {};
  \node[font=\Large] at (1,1.5) {\ja{カ}};
  \node[font=\scriptsize, text=gray!50!black] at (1,1.92) {ka};
  \node[cell] at (2,1.6) {};
  \node[font=\Large] at (2,1.5) {\ja{キ}};
  \node[font=\scriptsize, text=gray!50!black] at (2,1.92) {ki};
  \node[cell] at (3,1.6) {};
  \node[font=\Large] at (3,1.5) {\ja{ク}};
  \node[font=\scriptsize, text=gray!50!black] at (3,1.92) {ku};
  \node[cell] at (4,1.6) {};
  \node[font=\Large] at (4,1.5) {\ja{ケ}};
  \node[font=\scriptsize, text=gray!50!black] at (4,1.92) {ke};
  \node[cell] at (5,1.6) {};
  \node[font=\Large] at (5,1.5) {\ja{コ}};
  \node[font=\scriptsize, text=gray!50!black] at (5,1.92) {ko};
  \node[font=\small\bfseries, anchor=east] at (0.35,2.6) {s};
  \node[cell] at (1,2.6) {};
  \node[font=\Large] at (1,2.5) {\ja{サ}};
  \node[font=\scriptsize, text=gray!50!black] at (1,2.92) {sa};
  \node[cell] at (2,2.6) {};
  \node[font=\Large] at (2,2.5) {\ja{シ}};
  \node[font=\scriptsize, text=gray!50!black] at (2,2.92) {shi};
  \node[cell] at (3,2.6) {};
  \node[font=\Large] at (3,2.5) {\ja{ス}};
  \node[font=\scriptsize, text=gray!50!black] at (3,2.92) {su};
  \node[cell] at (4,2.6) {};
  \node[font=\Large] at (4,2.5) {\ja{セ}};
  \node[font=\scriptsize, text=gray!50!black] at (4,2.92) {se};
  \node[cell] at (5,2.6) {};
  \node[font=\Large] at (5,2.5) {\ja{ソ}};
  \node[font=\scriptsize, text=gray!50!black] at (5,2.92) {so};
  \node[font=\small\bfseries, anchor=east] at (0.35,3.6) {t};
  \node[cell] at (1,3.6) {};
  \node[font=\Large] at (1,3.5) {\ja{タ}};
  \node[font=\scriptsize, text=gray!50!black] at (1,3.92) {ta};
  \node[cell] at (2,3.6) {};
  \node[font=\Large] at (2,3.5) {\ja{チ}};
  \node[font=\scriptsize, text=gray!50!black] at (2,3.92) {chi};
  \node[cell] at (3,3.6) {};
  \node[font=\Large] at (3,3.5) {\ja{ツ}};
  \node[font=\scriptsize, text=gray!50!black] at (3,3.92) {tsu};
  \node[cell] at (4,3.6) {};
  \node[font=\Large] at (4,3.5) {\ja{テ}};
  \node[font=\scriptsize, text=gray!50!black] at (4,3.92) {te};
  \node[cell] at (5,3.6) {};
  \node[font=\Large] at (5,3.5) {\ja{ト}};
  \node[font=\scriptsize, text=gray!50!black] at (5,3.92) {to};
  \node[font=\small\bfseries, anchor=east] at (0.35,4.6) {n};
  \node[cell] at (1,4.6) {};
  \node[font=\Large] at (1,4.5) {\ja{ナ}};
  \node[font=\scriptsize, text=gray!50!black] at (1,4.92) {na};
  \node[cell] at (2,4.6) {};
  \node[font=\Large] at (2,4.5) {\ja{ニ}};
  \node[font=\scriptsize, text=gray!50!black] at (2,4.92) {ni};
  \node[cell] at (3,4.6) {};
  \node[font=\Large] at (3,4.5) {\ja{ヌ}};
  \node[font=\scriptsize, text=gray!50!black] at (3,4.92) {nu};
  \node[cell] at (4,4.6) {};
  \node[font=\Large] at (4,4.5) {\ja{ネ}};
  \node[font=\scriptsize, text=gray!50!black] at (4,4.92) {ne};
  \node[cell] at (5,4.6) {};
  \node[font=\Large] at (5,4.5) {\ja{ノ}};
  \node[font=\scriptsize, text=gray!50!black] at (5,4.92) {no};
  \node[font=\small\bfseries, anchor=east] at (0.35,5.6) {h};
  \node[cell] at (1,5.6) {};
  \node[font=\Large] at (1,5.5) {\ja{ハ}};
  \node[font=\scriptsize, text=gray!50!black] at (1,5.92) {ha};
  \node[cell] at (2,5.6) {};
  \node[font=\Large] at (2,5.5) {\ja{ヒ}};
  \node[font=\scriptsize, text=gray!50!black] at (2,5.92) {hi};
  \node[cell] at (3,5.6) {};
  \node[font=\Large] at (3,5.5) {\ja{フ}};
  \node[font=\scriptsize, text=gray!50!black] at (3,5.92) {fu};
  \node[cell] at (4,5.6) {};
  \node[font=\Large] at (4,5.5) {\ja{ヘ}};
  \node[font=\scriptsize, text=gray!50!black] at (4,5.92) {he};
  \node[cell] at (5,5.6) {};
  \node[font=\Large] at (5,5.5) {\ja{ホ}};
  \node[font=\scriptsize, text=gray!50!black] at (5,5.92) {ho};
  \node[font=\small\bfseries, anchor=east] at (0.35,6.6) {m};
  \node[cell] at (1,6.6) {};
  \node[font=\Large] at (1,6.5) {\ja{マ}};
  \node[font=\scriptsize, text=gray!50!black] at (1,6.92) {ma};
  \node[cell] at (2,6.6) {};
  \node[font=\Large] at (2,6.5) {\ja{ミ}};
  \node[font=\scriptsize, text=gray!50!black] at (2,6.92) {mi};
  \node[cell] at (3,6.6) {};
  \node[font=\Large] at (3,6.5) {\ja{ム}};
  \node[font=\scriptsize, text=gray!50!black] at (3,6.92) {mu};
  \node[cell] at (4,6.6) {};
  \node[font=\Large] at (4,6.5) {\ja{メ}};
  \node[font=\scriptsize, text=gray!50!black] at (4,6.92) {me};
  \node[cell] at (5,6.6) {};
  \node[font=\Large] at (5,6.5) {\ja{モ}};
  \node[font=\scriptsize, text=gray!50!black] at (5,6.92) {mo};
  \node[font=\small\bfseries, anchor=east] at (0.35,7.6) {y};
  \node[cell] at (1,7.6) {};
  \node[font=\Large] at (1,7.5) {\ja{ヤ}};
  \node[font=\scriptsize, text=gray!50!black] at (1,7.92) {ya};
  \node[cell, fill=gray!15] at (2,7.6) {};
  \node[cell] at (3,7.6) {};
  \node[font=\Large] at (3,7.5) {\ja{ユ}};
  \node[font=\scriptsize, text=gray!50!black] at (3,7.92) {yu};
  \node[cell, fill=gray!15] at (4,7.6) {};
  \node[cell] at (5,7.6) {};
  \node[font=\Large] at (5,7.5) {\ja{ヨ}};
  \node[font=\scriptsize, text=gray!50!black] at (5,7.92) {yo};
  \node[font=\small\bfseries, anchor=east] at (0.35,8.6) {r};
  \node[cell] at (1,8.6) {};
  \node[font=\Large] at (1,8.5) {\ja{ラ}};
  \node[font=\scriptsize, text=gray!50!black] at (1,8.92) {ra};
  \node[cell] at (2,8.6) {};
  \node[font=\Large] at (2,8.5) {\ja{リ}};
  \node[font=\scriptsize, text=gray!50!black] at (2,8.92) {ri};
  \node[cell] at (3,8.6) {};
  \node[font=\Large] at (3,8.5) {\ja{ル}};
  \node[font=\scriptsize, text=gray!50!black] at (3,8.92) {ru};
  \node[cell] at (4,8.6) {};
  \node[font=\Large] at (4,8.5) {\ja{レ}};
  \node[font=\scriptsize, text=gray!50!black] at (4,8.92) {re};
  \node[cell] at (5,8.6) {};
  \node[font=\Large] at (5,8.5) {\ja{ロ}};
  \node[font=\scriptsize, text=gray!50!black] at (5,8.92) {ro};
  \node[font=\small\bfseries, anchor=east] at (0.35,9.6) {w};
  \node[cell] at (1,9.6) {};
  \node[font=\Large] at (1,9.5) {\ja{ワ}};
  \node[font=\scriptsize, text=gray!50!black] at (1,9.92) {wa};
  \node[cell, fill=gray!15] at (2,9.6) {};
  \node[cell, fill=gray!15] at (3,9.6) {};
  \node[cell, fill=gray!15] at (4,9.6) {};
  \node[cell] at (5,9.6) {};
  \node[font=\Large] at (5,9.5) {\ja{ヲ}};
  \node[font=\scriptsize, text=gray!50!black] at (5,9.92) {o};
  \node[font=\small\bfseries, anchor=east] at (0.35,10.6) {};
  \node[cell] at (1,10.6) {};
  \node[font=\Large] at (1,10.5) {\ja{ン}};
  \node[font=\scriptsize, text=gray!50!black] at (1,10.92) {n};
  \node[cell, fill=gray!15] at (2,10.6) {};
  \node[cell, fill=gray!15] at (3,10.6) {};
  \node[cell, fill=gray!15] at (4,10.6) {};
  \node[cell, fill=gray!15] at (5,10.6) {};
\end{tikzpicture}
\end{center}

Dakuten, handakuten, small \ja{ャュョ} and small \ja{ッ} work exactly as in hiragana:
\ja{ガ} ga, \ja{パ} pa, \ja{キャ} kya, \ja{ベッド} \emph{beddo} ``bed''.

\subsection{Where the shapes come from}

Each katakana is a fragment of a kanji with (roughly) that sound. Knowing the source
helps a few shapes stick:
\begin{center}
\begin{tikzpicture}[src/.style={draw, minimum size=1.2cm, font=\LARGE},
    kana/.style={draw, fill=blue!8, minimum size=1.2cm, font=\LARGE}]
  \foreach \s/\k/\r [count=\i] in {加/カ/ka, 伊/イ/i, 多/タ/ta, 千/チ/chi, 八/ハ/ha, 江/エ/e} {
    \node[src] (s\i) at (\i*2.2, 0) {\ja{\s}};
    \node[kana] (k\i) at (\i*2.2, -2) {\ja{\k}};
    \draw[-Stealth, thick] (s\i) -- (k\i);
    \node[font=\small] at (\i*2.2, -2.95) {\r};
  }
  \node[anchor=east, font=\small] at (1.4,0) {kanji};
  \node[anchor=east, font=\small] at (1.4,-2) {katakana};
\end{tikzpicture}
\end{center}
\ja{カ} is the left part of \ja{加}, \ja{イ} the left radical of \ja{伊}, \ja{タ} the top of
\ja{多}, \ja{エ} the right side of \ja{江}; \ja{チ} and \ja{ハ} are simplified whole kanji.

\subsection{Confusable katakana}

\begin{center}
\begin{tabular}{ll}
\toprule
Pair & How to tell them apart \\
\midrule
\ja{シ} shi / \ja{ツ} tsu & \ja{シ}: dots stacked on the left, last stroke sweeps \emph{up}; \ja{ツ}: dots in a row, stroke sweeps \emph{down} \\
\ja{ソ} so / \ja{ン} n & \ja{ソ}: stroke falls from the top; \ja{ン}: stroke rises from the bottom \\
\ja{ク} ku / \ja{ワ} wa / \ja{ケ} ke & \ja{ワ} has a flat top; \ja{ケ} has a separate cross bar \\
\ja{ス} su / \ja{ヌ} nu & \ja{ヌ} has an extra short slash \\
\ja{コ} ko / \ja{ユ} yu & \ja{ユ}'s bottom line sticks out to the right \\
\ja{ウ} u / \ja{フ} fu & \ja{ウ} has a dot and a short vertical on top \\
\ja{チ} chi / \ja{テ} te & \ja{テ} has two flat bars and no sloping top \\
\bottomrule
\end{tabular}
\end{center}

\begin{center}
\begin{tikzpicture}[x=1cm,y=1cm]
  % shi
  \draw[thick] (0,0) rectangle (2,2);
  \draw[very thick, -Stealth] (0.35,1.6) -- (0.6,1.45);
  \draw[very thick, -Stealth] (0.3,1.05) -- (0.55,0.9);
  \draw[very thick, -Stealth] (0.4,0.2) .. controls (1,0.4) .. (1.65,1.5);
  \node at (1,-0.4) {\ja{シ} shi: dots stacked, stroke rises};
  % tsu
  \begin{scope}[xshift=5cm]
  \draw[thick] (0,0) rectangle (2,2);
  \draw[very thick, -Stealth] (0.3,1.6) -- (0.45,1.25);
  \draw[very thick, -Stealth] (0.85,1.65) -- (1,1.3);
  \draw[very thick, -Stealth] (1.7,1.65) .. controls (1.6,0.6) .. (0.7,0.2);
  \node at (1,-0.4) {\ja{ツ} tsu: dots side by side, stroke falls};
  \end{scope}
\end{tikzpicture}
\end{center}

\section{Writing loanwords}

\subsection{The long-vowel bar}
In katakana a long vowel is written with a bar \ja{ー} (vertical in vertical writing)
rather than an extra vowel letter:
\ja{コーヒー} \emph{kōhī} coffee, \ja{スーパー} \emph{sūpā} supermarket,
\ja{ノート} \emph{nōto} notebook, \ja{ケーキ} \emph{kēki} cake.

\subsection{Extended combinations}
To approximate foreign sounds, katakana uses small vowels \ja{ァィゥェォ} in combinations
that do not exist in hiragana:
\begin{center}
\begin{tabular}{llll}
\toprule
Combo & Sound & Example & Meaning \\
\midrule
\ja{ファ フィ フェ フォ} & fa fi fe fo & \ja{フォーク} fōku & fork \\
\ja{ティ ディ} & ti di & \ja{パーティー} pātī & party \\
\ja{トゥ ドゥ} & tu du & \ja{タトゥー} tatū & tattoo \\
\ja{ウィ ウェ ウォ} & wi we wo & \ja{ウェブ} webu & web \\
\ja{シェ ジェ チェ} & she je che & \ja{チェック} chekku & check \\
\ja{ヴ} & vu (often just b) & \ja{ヴァイオリン} vaiorin & violin \\
\bottomrule
\end{tabular}
\end{center}

\subsection{How English words get ``Japanified''}
\begin{enumerate}
  \item Every consonant (except n) must be followed by a vowel. Usually \emph{u} is added;
        after t and d, \emph{o} is added: ``test'' $\to$ \ja{テスト} \emph{tesuto};
        ``bed'' $\to$ \ja{ベッド} \emph{beddo}.
  \item l and r both become the r-row: ``light'' and ``right'' $\to$ \ja{ライト}
        \emph{raito}.
  \item Stressed short vowels before a final consonant often trigger \ja{ッ}: ``cup''
        $\to$ \ja{カップ} \emph{kappu}.
  \item ``-er'', ``-ar'' endings become a long vowel: ``computer'' $\to$
        \ja{コンピューター} \emph{konpyūtā}.
  \item ``th'' becomes s or z: ``three'' $\to$ \ja{スリー} \emph{surī}.
\end{enumerate}

\begin{example}[Our names in katakana]
\leavevmode
\begin{itemize}
  \item Vishesh $\to$ \ja{ヴィシェシュ} \emph{Visheshu} (or \ja{ビシェシュ}).
  \item Emily $\to$ \ja{エミリー} \emph{Emirī}.
  \item Michael $\to$ \ja{マイケル} \emph{Maikeru}.
  \item Kevin $\to$ \ja{ケビン} \emph{Kebin}.
  \item Sarah $\to$ \ja{サラ} \emph{Sara}.
\end{itemize}
\end{example}

\section{Loanword vocabulary}

\begin{center}
\begin{tabular}{lll}
\toprule
Katakana & Romaji & Meaning (source) \\
\midrule
\ja{アルバイト} & arubaito & part-time job (German \emph{Arbeit}) \\
\ja{パン} & pan & bread (Portuguese \emph{pão}) \\
\ja{テレビ} & terebi & television \\
\ja{パソコン} & pasokon & PC (\emph{perso}nal \emph{com}puter) \\
\ja{スマホ} & sumaho & smartphone \\
\ja{コンビニ} & konbini & convenience store \\
\ja{レストラン} & resutoran & restaurant \\
\ja{アメリカ} & Amerika & America (USA) \\
\ja{イギリス} & Igirisu & the UK (Portuguese \emph{Inglês}) \\
\ja{ホッケー} & hokkē & hockey \\
\bottomrule
\end{tabular}
\end{center}

\begin{aside}
Many loanwords are clipped (\ja{パソコン}, \ja{スマホ}, \ja{コンビニ}) or have shifted
meaning (\ja{マンション} \emph{manshon} is a condo, not a mansion). These \emph{wasei
eigo} (``Japan-made English'') words are not understood by English speakers, so do not
assume a katakana word means what its English source means.
\end{aside}

\section*{Lesson 3 review}
\begin{itemize}
  \item Katakana: same sounds as hiragana, angular shapes from kanji fragments.
  \item Used for loanwords, foreign names, sound effects and emphasis.
  \item Long vowels use \ja{ー}; extra combinations (\ja{ファ}, \ja{ティ}, \ja{ウェ}, \ldots)
        approximate foreign sounds.
  \item Watch \ja{シ}/\ja{ツ} and \ja{ソ}/\ja{ン}: stroke direction is the difference.
\end{itemize}
